\documentclass[11pt]{article}
\usepackage[margin=1in]{geometry}
\usepackage{amsmath,amssymb,mathtools}
\usepackage{booktabs}
\usepackage{enumitem}
\usepackage{siunitx}
\sisetup{group-separator={,},group-minimum-digits=3}

\title{\textbf{Quiz 1} \vspace{ 0.4cm }\\ECON 300: Labour Economics}
\author{Instructor: Muhebullah Karimzada}
%\date{January 22,  2026}

\begin{document}
\maketitle

\section*{Part A: Multiple Choice (5 $\times$ 1 = 5 points)}

\textbf{Choose the best answer for each question.}

\begin{enumerate}
\item Which of the following individuals is counted as \textbf{not in the labour force}?

\begin{enumerate}
\item A worker who searched for a job last week  
\item A worker on temporary layoff  
\item A discouraged worker who stopped searching  
\item A worker who starts a new job next week  
\end{enumerate}

\item The employment-to-population ratio measures:

\begin{enumerate}
\item Employment divided by the labour force  
\item Employment divided by the working-age population  
\item The fraction of unemployed workers  
\item The fraction of discouraged workers  
\end{enumerate}

\item In the labour--leisure model, an increase in non-labour income:

\begin{enumerate}
\item Rotates the budget constraint  
\item Has no effect on labour supply  
\item Shifts the budget constraint upward in a parallel way  
\item Makes leisure cheaper relative to consumption  
\end{enumerate}

\item A backward-bending labour supply curve arises when:

\begin{enumerate}
\item The substitution effect dominates at all wages  
\item Leisure is an inferior good  
\item The income effect dominates at high wages  
\item Non-labour income is zero  
\end{enumerate}


\newpage
\item Which of the following policies is most likely to \textbf{reduce} measured labour force participation?

\begin{enumerate}
\item A rise in wages  
\item An expansion of unemployment insurance  
\item Improved job matching technology  
\item Lower payroll taxes  
\end{enumerate}
\end{enumerate}



\subsection*{Problem 1: Labour Force Measurement and Discouragement (7.5 points)}

A city has a working-age population (15+) of 900{,}000 people. The initial breakdown is:

\begin{itemize}
\item Employed: 585{,}000  
\item Unemployed: 65{,}000  
\item Not in the labour force: 250{,}000  
\end{itemize}

\begin{enumerate}
\item[(a)] (2.5 points)  
Calculate the \textbf{labour force}, the \textbf{labour force participation rate}, and the \textbf{unemployment rate}.

\item[(b)] (2.5 points)  
Calculate the \textbf{employment-to-population ratio} and explain how it differs conceptually from the LFPR.

\item[(c)] (2.5 points)  
Suppose an economic slowdown causes:
\begin{itemize}
\item 15{,}000 employed workers to lose their jobs, and  
\item 10{,}000 unemployed workers to stop searching for work.
\end{itemize}

Recalculate the unemployment rate and explain why changes in discouragement complicate the interpretation of labour market conditions.
\end{enumerate}

\subsection*{Problem 2: Labour Supply and Minimum Wage (7.5 points)}

A worker's weekly budget constraint is:

\[
C = 500 + 30h
\]

and preferences are given by:

\[
U = C - h^2
\]

The labour market is described by:

\[
L^D = 150 - 3w,
\qquad
L^S = 30 + 2w
\]

\begin{enumerate}
\item[(a)] (2.5 points)  
Find the worker's \textbf{optimal hours of work}.

\item[(b)] (2.5 points)  
Find the \textbf{competitive equilibrium wage and employment}.

\item[(c)] (2.5 points)  
A minimum wage of \( w = 25 \) is introduced.  
Determine whether it is binding and calculate the resulting level of \textbf{unemployment}.
\end{enumerate}

\vfill
\textbf{End of Quiz}

\end{document}
